\documentclass[aspectratio = 1610, french]{beamer}

\usetheme{Madrid}

\setbeamercovered{transparent}
\setbeamertemplate{caption}[numbered]

% Packages

%% Font configuration

\usepackage[T1]{fontenc}
\usepackage{lmodern}

%% i18n packages

\usepackage[main = french, english]{babel}
\usepackage{fvextra} % needed to be loaded before `csquotes`
\usepackage[autostyle = true]{csquotes} % for `\enquote` command
\usepackage[useregional]{datetime2} % better date formatting
\usepackage{icomma} % fix comma space

%% Main packages

\usepackage[french]{cleveref}

\usepackage[
    type = {CC},
    modifier = {by-nc},
    version = {4.0},
]{doclicense}

\usepackage{nicefrac}
\usepackage{physics}
\usepackage{siunitx}
\usepackage[most]{tcolorbox}
\usepackage{wrapfig}


%% `tabular` packages

\usepackage{multicol}
\usepackage{multirow}
\usepackage{tabularx}

%% Drawing packages

\usepackage{svg}
\usepackage{tikz}

% Some command overrides

\renewcommand{\arraystretch}{1.5}

%% Patch cleveref

%%% See https://tex.stackexchange.com/a/670393.

\makeatletter
\patchcmd\@cref{\begingroup}{\begingroup\@safe@activestrue}{}{\fail}
\patchcmd\@cref{\cref@isstackfull{\@refsubstack}}{\cref@isstackfull{\@refsubstack}\@safe@activesfalse}{}{\fail}
\makeatother

%% Section headers

\renewcommand{\thesection}{\Roman{section}}
\renewcommand{\thesubsection}{\arabic{subsection}}
\renewcommand{\thesubsubsection}{\alph{subsubsection}}

% New commands

\DeclareMathOperator{\sign}{sign}

%% Text commands

\newcommand{\textdef}[1]{\textcolor{red}{\text{#1}}}
\newcommand{\textimp}[1]{\textcolor{teal}{\text{#1}}}

% `tabularx`

\newcolumntype{Y}{>{\centering\arraybackslash}X}

% Some `tcolorbox` definitions

%% Styles

\tcbset{
    tab/.style = {
        center title,
        colback = white!10!white,
        colbacktitle = Salmon!40!white,
        colframe = red!50!black,
        coltitle = black,
        enhanced,
        fonttitle = \bfseries,
        fontupper = \normalsize \sffamily
    }
}

%% Environments

\newtcolorbox{tcbtabular}[2]{
    tab,
    tabularx = {#1},
    #2
}

% General setup

\author[]{
    Otman BENJELLOUN (\href{https://github.com/otman-benjelloun/}{@otman-benjelloun}) \and
    Emilia FER (\href{https://github.com/emilia-fer/}{@emilia-fer}) \and
    Eline GEORGE (\href{https://github.com/elinegeorge/}{@elinegeorge}) \and
    Mattéo ROSSILLOL-{}-LARUELLE (\href{https://github.com/beatussum/}{@beatussum}) \and
    \href{https://github.com/tvidon/}{@tvidon}
}

\title[4MMPSPE -- RLST]{
    Projet de spécialité -- RLST \\
    \small{4MMPSPE}
}

\begin{document}

\frame{\titlepage}

\begin{frame}

\frametitle{Sommaire}
\tableofcontents[pausesections, pausesubsections]

\end{frame}

\section{Avant-propos}

\begin{frame}

\frametitle{Avant-propos}

Ce projet spécialité portait sur la réalisation d'un outil de synthèse logique permettant la compilation d'un code Verilog en circuit \emph{Redstone} Minecraft. De manière plus factuelle, l'objectif était donc

\begin{enumerate}
    \item<+-> de lancer YOSYS à partir d'un fichier Verilog,
    \item<+-> associer les portes avec leur représentation dans \emph{Minecraft},
    \item<+-> de placer les portes logiques récupérées,
    \item<+-> de router les portes, et
    \item<+-> de les placer dans un fichier MCA.
\end{enumerate}

\end{frame}

\section{Structure du projet}

\begin{frame}

\frametitle{Structure du projet}

\begin{actionenv}<+->

Le projet se décompose en deux parties principales :

\begin{itemize}
    \item<+-> \texttt{rlst\_core}, une bibliothèque statique qui implémente l'essentielle de l'\emph{API} ;
    \item<+-> \texttt{rlst}, un exécutable servant de lanceur.
\end{itemize}

\end{actionenv}

\begin{actionenv}<+->

La bibliothèque est plusieurs en trois \emph{packages} :

\begin{itemize}
    \item<+-> \texttt{cell\_reader} : lecture des fichiers décrivant les portes logiques ;
    \item<+-> \texttt{par} (\foreignlanguage{english}{Place And Route}) ;
    \item<+-> \texttt{svg} : utilitaire permettant de faire du dessin vectoriel (utilisé pour le débogage).
\end{itemize}

\end{actionenv}

\end{frame}

\subsection{Lecture du VHDL avec YOSYS}

\begin{frame}

\frametitle{Les portes logiques}

Les portes logiques se caractérisent par

\begin{itemize}
    \item<+-> leur taille,
    \item<+-> la position relative des ports au sein de la portes,
    \item<+-> les blocs qui la composent.
\end{itemize}

\end{frame}

\subsection{Placement}

\begin{frame}

\frametitle{Recuit simulé}

L'algorithme de placement repose sur du recuit simulé.

\begin{actionenv}<+->

Les coûts évaluent

\begin{itemize}
    \item<+-> la proportion des chevauchements (\textimp{bornée naturellement}),
    \item<+-> la longueur estimée des fils (\textimp{normalisée}),
    \item<+-> la proportion de cases vides (\textimp{bornée naturellement}),
    \item<+-> l'écart relatif séparant la moyenne de la longueur des lignes de celle désirée (\textimp{normalisée}),
    \item<+-> la hauteur de la configuration (\textimp{normalisée}), et
    \item<+-> l'espace séparant l'axe des abscisses de la cellule la plus haute (\textimp{normalisée}).
\end{itemize}

\end{actionenv}

\end{frame}

\begin{frame}

\frametitle{Fonction d'acceptation et température initiale}

La \textimp{fonction d'acceptation} et la \textimp{température initiale} sont de

\begin{align}
\exp \biggl( \frac{100 \Delta C}{T} \biggr) && T_0 = 500
\end{align}

Empiriquement, on a montré que \( \Delta C \simeq 1 \) et \( T \simeq 100 \) ; ainsi pour que ces deux grandeurs est le même ordre de grandeur on multiplie \( \Delta C \) par un facteur 100.

\end{frame}

\begin{frame}

\frametitle{Fonction de normalisation}

Afin d'avoir des coûts avec des valeurs bornées dans \( [0 ; 100] \), il a fallu utilisé une \textimp{fonction de normalisation} :

\begin{align}
\nicefrac{100}{\pi} \cdot \sign x \cdot \atan \Bigl( \log \abs{x} + 1 \Bigr) + 50
\end{align}

Cette fonction a le bon goût

\begin{itemize}
    \item<+-> d'être bornée,
    \item<+-> symétrique,
    \item<+-> d'avoir une croissance relativement faible, et
    \item<+-> coupe l'axe des ordonnées à 50 en \( x = 0 \).
\end{itemize}

\end{frame}

\subsection{Routage}

\begin{frame}

\frametitle{Principe de base}

\begin{actionenv}<+->
Pour router les cellules, on s'appuie sur le fait que \textimp{Minecraft} permette d'établir des structures \emph{électrique} en jouant sur la \textimp{verticalité}.
\end{actionenv}

\begin{actionenv}<+->
Ainsi, les routes se situent à plusieurs niveaux différents de telle manière à ce que les tracés d'un même niveau ne s'intersecte pas.
\end{actionenv}

\begin{actionenv}<+->
Cependant, on note qu'il existe une petite subtilité afin d'éviter que les routes viennent interagir avec des portes alors qu'elles ne devraient pas : \textimp{on dilate virtuellement les portes}.
\end{actionenv}

\end{frame}

\begin{frame}

\frametitle{Point sur les répéteurs}

Un signal de \emph{Redstone} sous Minecraft est caractérisé par

\begin{itemize}
    \item<+-> une puissance variant entre 0 et 15,
    \item<+-> cette puissance se dissipe d'une unité pour chaque bloc parcouru,
    \item<+-> la puissance peut être re-maximiser avec un bloc particulier, le \emph{repeter}.
\end{itemize}

\end{frame}

\subsection{Édition du monde Minecraft}

\begin{frame}

\frametitle{Format du monde}

Un \textimp{monde \emph{Minecraft}} est divisé en plusieurs

\begin{itemize}
    \item<+-> \textimp{régions} de \( 32 \times 32 = \num{1024} \) \textimp{chunks},
    \item<+-> une portion de région de 16 blocs de longueur, 16 blocs de largeur et 256 blocs de hauteur.
\end{itemize}

\begin{actionenv}<+->

Pour utiliser ce format, il a fallu veiller

\begin{itemize}
    \item<+-> au boutisme, et
    \item<+-> au format de compression (\textimp{ZLIB}) des \emph{chunks}.
\end{itemize}

\end{actionenv}

\end{frame}

\section{Synthèse}

\begin{frame}

\frametitle{Synthèse}

Ainsi, on a pu implémenter un programme permettant

\begin{enumerate}
    \item<+-> de lancer YOSYS à partir d'un fichier Verilog (\textimp{bogué}),
    \item<+-> associer les portes avec leur représentation dans \emph{Minecraft},
    \item<+-> de placer les portes logiques récupérées,
    \item<+-> de router les portes, et
    \item<+-> de les placer dans un fichier MCA (\textimp{bogué}).
\end{enumerate}

\begin{actionenv}<+->
Bien que le projet ne soit pas parfaitement fini, on a pu montrer que l'objectif visé était réalisable, et on a pu montrer une première esquisse du résultat attendu.
\end{actionenv}

\end{frame}

\begin{frame}

\begin{figure}[t]
\centering\Large
Merci pour votre écoute.
\end{figure}

\begin{figure}[b]
\doclicenseThis
\end{figure}

\end{frame}

\end{document}
